\documentclass[oneside,a4paper,14pt]{extarticle}
\usepackage[a4paper,top=20mm,bottom=20mm,left=30mm,right=10mm]{geometry}
\usepackage{amsmath, amssymb, amsfonts, amsthm}
\usepackage[T2A]{fontenc}
\usepackage[utf8]{inputenc}
\usepackage[russian]{babel}
\usepackage{textcomp}
\usepackage{indentfirst}
\usepackage{graphicx}
\usepackage{float}
\usepackage{wrapfig}
\usepackage{caption}
\usepackage{tocloft}
\renewcommand{\cftsecleader}{\cftdotfill{\cftdotsep}}
\usepackage[breaklinks]{hyperref}
\usepackage{titlesec}
\usepackage{fancyvrb}
\usepackage{xcolor}
\usepackage{listings}

\titleformat{\section}{\normalsize\bfseries}{\thesection}{1em}{}
\titleformat{\subsection}{\normalsize\bfseries}{\thesubsection}{1em}{}
\titleformat{\subsubsection}{\normalsize\bfseries}{\thesubsubsection}{1em}{}
\renewcommand\baselinestretch{1.33}\normalsize
\setlength{\parindent}{1.25cm}
\begin{document}
\newpage\thispagestyle{empty}
\newpage
\setcounter{page}{1}
\tableofcontents
\newpage

\section*{Тема}
\addcontentsline{toc}{section}{Тема}
«Получение прогноза погоды из OpenWeatherMap, определение местоположения через IP Geolocation API, сохранение в SQLite, выгрузка в Markdown и параметризация подключения к БД через переменные окружения»

\section*{Цель работы}
\addcontentsline{toc}{section}{Цель работы}
Освоить базовый ETL-контур на Python:
\begin{itemize}
	\item автоматическое определение местоположения через IP-геолокацию;
	\item получение прогноза погоды на сегодня и ближайшие 3 дня из OpenWeatherMap;
	\item нормализацию и агрегацию данных;
	\item сохранение в локальной СУБД (SQLite);
	\item выгрузку обработанных данных в файл формата Markdown с красивым оформлением;
	\item управление параметрами подключения к базе данных через переменные окружения;
	\item публикацию проекта на GitHub с правильной организацией репозитория.
\end{itemize}

\section*{Используемые технологии и API}
\addcontentsline{toc}{section}{Используемые технологии и API}
Для выполнения лабораторной работы были использованы следующие технологии и API:
\begin{itemize}
	\item Python - основной язык разработки; реализованы сбор данных, парсинг, работа с БД и генерация отчёта.
	\item requests - HTTP-клиент для обращения к внешним REST-API (OpenWeatherMap и ipinfo.io), поддерживает параметры запроса, таймауты и обработку исключений.
	\item SQLAlchemy - ORM для описания модели WeatherPoint, создания таблиц и выполнения запросов select/insert без написания SQL вручную.
	\item python-dotenv - загрузка переменных окружения из файла .env, что позволяет не хранить секреты (API-ключ, строку подключения) в репозитории.
	\item SQLite - встраиваемая СУБД; файл weather.db создаётся автоматически, не требует отдельного сервера.
	\item OpenWeatherMap - источник прогноза погоды; используется endpoint /forecast с параметрами lat, lon, units=metric, lang=ru.
	\item ipinfo.io - сервис IP-геолокации; по запросу https://ipinfo.io/json возвращает координаты и город, которые затем передаются в OpenWeatherMap.
\end{itemize}

\section*{Структура проекта}
\addcontentsline{toc}{section}{Структура проекта}
Исходя из требований лабораторной работы, структура проекта в gitHub выглядит так:

\begin{small}
	\begin{verbatim}
Vyz/3курс/СЦИФУПР/1/weather-etl-lab/
- .env               <- не в репозитории
- .env.example       <- шаблон
- README.md
- requirements.txt
- main.py
- api_caller.py
- config.py
- save.py
- to_md.py
- wpp_parser.py
\end{verbatim}
\end{small}

\section*{Ключевые фрагменты кода}
\addcontentsline{toc}{section}{Ключевые фрагменты кода}

\subsection*{Конфигурация}
\addcontentsline{toc}{subsection}{Конфигурация}
Все параметры подключения вынесены в переменные окружения. Значения по умолчанию заданы для url, чтобы код работал даже при неполном .env.

\begin{small}
	\begin{verbatim}
import os
from dotenv import load_dotenv

load_dotenv()

WEATHER_API_KEY: str = os.getenv("WEATHER_API_KEY")
WEATHER_API_BASE_URL: str = os.getenv(
    "WEATHER_API_BASE_URL", "https://api.openweathermap.org/data/2.5"
)
GEO_API_URL: str = os.getenv("GEO_API_URL")
GEO_API_FALLBACK_CITY: str = os.getenv("GEO_API_FALLBACK_CITY")
DB_URL: str = os.getenv("DB_URL")
\end{verbatim}
\end{small}

\subsection*{HTTP-клиент}
\addcontentsline{toc}{subsection}{HTTP-клиент}
Функция get\_data - единая точка обращения к внешним API. Она выполняет запрос, проверяет статус и логирует специфические ошибки.

\begin{small}
	\begin{verbatim}
def get_data(API_URL: str, parametrs: dict[str, Any] | None = None) -> dict[str, Any]:
    try:
        response = requests.get(url=API_URL, params=parametrs)
        response.raise_for_status()
        return response.json()
    except requests.exceptions.Timeout:
        logger.error("Таймаут при запросе к %s", API_URL)
        raise
    except requests.exceptions.ConnectionError:
        logger.error("Ошибка соединения с %s", API_URL)
        raise
    except requests.exceptions.HTTPError as exc:
        status = exc.response.status_code
        if status == 401:
            logger.error("Неверный API-ключ.")
        elif status == 404:
            logger.error("Ресурс не найден.")
        elif status == 429:
            logger.error("Превышен лимит запросов.")
        else:
            logger.error("Ошибка %s: %s", status, exc)
        raise
    except requests.exceptions.RequestException as exc:
        logger.error("Неизвестная ошибка запроса: %s", exc)
        raise
\end{verbatim}
\end{small}

\subsection*{Парсинг прогноза}
\addcontentsline{toc}{subsection}{Парсинг прогноза}
OpenWeatherMap возвращает точки с шагом 3 часа. Из них отбираются записи за сегодня и следующие 3 дня, которые преобразуются в объекты WeatherPoint.

\begin{small}
	\begin{verbatim}
def parse_weather(raw_data: dict[str, Any]) -> list[WeatherPoint]:
    DAYS: list[WeatherPoint] = []
    list_of_points: list[dict[str, Any]] = raw_data["list"]

    today = datetime.now().date()
    end_date = today + timedelta(days=3)

    for day in list_of_points:
        point_date = datetime.fromisoformat(day["dt_txt"]).date()
        if point_date < today or point_date > end_date:
            continue
        DAYS.append(WeatherPoint(...))
    return DAYS
\end{verbatim}
\end{small}

\subsection*{Сохранение в SQLite}
\addcontentsline{toc}{subsection}{Сохранение в SQLite}
Модель WeatherPoint описана через SQLAlchemy 2.0. Перед записью проверяются уже существующие даты, чтобы не дублировать записи.

\begin{small}
	\begin{verbatim}
class WeatherPoint(Base):
    __tablename__ = "weather"
    id: Mapped[int] = mapped_column(primary_key=True)
    city: Mapped[str] = mapped_column(String(100), default="")
    date: Mapped[date] = mapped_column(Date, index=True)
    temperature: Mapped[float] = mapped_column(Float, default=0.0)
    # ... остальные поля ...

def save_points(points: list[WeatherPointModel], city: str = "") -> None:
    with Session(engine) as session:
        existing_dates: set[date] = set(
            session.scalars(select(WeatherPoint.date)).all()
        )
        new_points = [p for p in points if p.date not in existing_dates]
        session.add_all(WeatherPoint(...) for p in new_points)
        session.commit()
\end{verbatim}
\end{small}

\subsection*{Выгрузка в Markdown}
\addcontentsline{toc}{subsection}{Выгрузка в Markdown}
Формируется таблица с датой, минимальной и максимальной температурой, описанием, влажностью и скоростью ветра.

\begin{small}
	\begin{verbatim}
lines: list[str] = [
    "# Прогноз погоды",
    "",
    f"Автоматически определённая локация: {city or 'не определена'}  ",
    f"Период: {period}  ",
    "",
    "| Дата | Мин. темп. | Макс. темп. | Описание | Влажность | Ветер |",
    "|------|-----------------|------------------|----------|---------------|-------------|",
]
for d in points:
    lines.append(
        f"| {date_str} | {d.temp_min:.1f} | {d.temp_max:.1f} "
        f"| {d.description} | {d.humidity:.0f} | {d.windspeed:.1f} |"
    )
\end{verbatim}
\end{small}

\section*{Обработка ошибок}
\addcontentsline{toc}{section}{Обработка ошибок}
Обработка ошибок построена по принципу «не падать молча»: каждая проблема логируется и, если она критична для дальнейших шагов, пробрасывается наверх.

\begin{itemize}
	\item Сетевые ошибки. В get\_data отдельно перехватываются Timeout, ConnectionError и общий RequestException. Для HTTP-ошибок различаются коды 401 (неверный API-ключ), 404 (ресурс не найден) и 429 (превышен лимит запросов).
	\item Ошибки формата ответа. В get\_geo проверяется наличие поля loc; если его нет - возбуждается KeyError с сообщением «Неправильный ответ от API».
	\item Ошибки конфигурации. В main проверяется, что WEATHER\_API\_KEY задан; иначе работа прекращается с сообщением в лог.
	\item Ошибки БД. Операции init\_db, save\_points, load\_points обёрнуты в try/except на уровне main, чтобы при сбое не терять уже собранные данные и не завершать программу аварийно.
	\item Логирование. Используется стандартный модуль logging с уровнем INFO и форматом, включающим время, уровень, имя логгера и сообщение. Для отладочных сведений (например, список точек) применяется logger.debug.
\end{itemize}

\section*{Результаты работы}
\addcontentsline{toc}{section}{Результаты работы}
В результате выполнения main.py:
\begin{enumerate}
	\item создаётся (или пересоздаётся) файл базы данных weather.db и таблица weather;
	\item через https://ipinfo.io/json определяется текущий IP, затем через IP-геолокацию - координаты и город;
	\item выполняется запрос к OpenWeatherMap /forecast и из ответа отбираются точки за сегодня и ближайшие 3 дня;
	\item отобранные точки сохраняются в SQLite (без дублей по дате);
	\item формируется файл weather\_report.md с заголовком, периодом, названием локации и таблицей прогноза;
	\item все ключевые шаги сопровождаются записями в логе.
\end{enumerate}

Пример содержимого weather\_report.md:

\begin{small}
	\begin{verbatim}
# Прогноз погоды

Автоматически определённая локация: Moscow
Период: 2025-01-01 - 2025-01-04

| Дата | Мин. темп. | Макс. темп. | Описание | Влажность | Ветер |
|------|-----------------|------------------|----------|---------------|-------------|
| 2025-01-01 | -3.5 | -1.0 | небольшой снег | 85 | 4.2 |
| 2025-01-02 | -4.1 | -2.3 | облачно | 80 | 3.1 |
\end{verbatim}
\end{small}

\section*{Содержимое .env.example}
\addcontentsline{toc}{section}{Содержимое .env.example}
Шаблонный файл окружения содержит следующие переменные:

\begin{small}
	\begin{verbatim}
WEATHER_API_KEY=apykey
WEATHER_API_BASE_URL=https://api.openweathermap.org/data/2.5

GEO_API_URL=https://ipinfo.io/json
GEO_API_FALLBACK_CITY=Moscow

DB_URL=sqlite:///weather.db

ENVIRONMENT=development
OUTPUT_FILE=weather_report.md
\end{verbatim}
\end{small}

\section*{Вывод}
\addcontentsline{toc}{section}{Вывод}
В ходе лабораторной работы был реализован полный ETL-контур: получение данных из внешних API, их нормализация, сохранение в реляционную БД и выгрузка в человекочитаемый формат Markdown. Освоены:
\begin{itemize}
	\item работа с REST-API через requests и обработка сетевых/HTTP-ошибок;
	\item параметризация конфигурации через python-dotenv и переменные окружения;
	\item описание модели данных и запросы к SQLite через SQLAlchemy 2.0;
	\item логирование с разделением уровней INFO/DEBUG;
	\item генерация отчёта в формате Markdown.
\end{itemize}
Полученный код легко расширяется: достаточно добавить новые поля в модель и парсер, чтобы включить в отчёт дополнительные метрики (например, давление или облачность).

\section*{Список литературы}
\addcontentsline{toc}{section}{Список литературы}
\begin{enumerate}
	\item OpenWeatherMap API - документация сервиса прогноза погоды. URL: \href{https://openweathermap.org/api}{https://openweathermap.org/api}
	\item IP Geolocation API (ipinfo.io) - определение местоположения по IP-адресу. URL: \href{https://ipinfo.io/developers}{https://ipinfo.io/developers}
	\item SQLite Documentation - встраиваемая СУБД. URL: \href{https://docs.python.org/3/library/sqlite3.html}{https://docs.python.org/3/library/sqlite3.html}
	\item SQLAlchemy Documentation - ORM и работа с БД на Python. URL: \href{https://docs.sqlalchemy.org/en/14/}{https://docs.sqlalchemy.org/en/14/}
	\item python-dotenv Documentation - загрузка переменных окружения из .env. URL: \href{https://pypi.org/project/python-dotenv/}{https://pypi.org/project/python-dotenv/}
	\item Python JSON Module Documentation - сериализация и десериализация JSON. URL: \href{https://docs.python.org/3/library/json.html}{https://docs.python.org/3/library/json.html}
	\item Python OS Module Documentation - работа с переменными окружения и файловой системой. URL: \href{https://docs.python.org/3/library/os.html}{https://docs.python.org/3/library/os.html}
	\item Python Logging Module Documentation - логирование в Python. URL: \href{https://docs.python.org/3/library/logging.html}{https://docs.python.org/3/library/logging.html}
	\item Python Unittest Module Documentation - модульное тестирование. URL: \href{https://docs.python.org/3/library/unittest.html}{https://docs.python.org/3/library/unittest.html}
	\item GitHub REST API Documentation - работа с репозиториями через API. URL: \href{https://docs.github.com/en/rest}{https://docs.github.com/en/rest}
\end{enumerate}
\end{document}